\begin{claim}{C.76}{Proved}{Simultaneous endpoint prescriptions with compact original-boundary control}
Let $Q$ be the actual closed, locally compact, free proper
one-dimensional group on $O=\operatorname{int}\overline D$
of C.75. Let $f_\tau$ be a based loop of homeomorphisms of
the actual closed disk, fixing its boundary pointwise and
commuting with $Q$ on $O$. Use the actual principal
coordinate and its inverse from C.75, with coordinates
$(s,t)$, common base fixed set $A$, and generated group $G$.

Retain the following components $I$ of $\mathbb R\setminus A$:
those for which every finite-endpoint character $c_a$ is zero,
and those having a finite endpoint $a$ with
\[
       c_a\ne0,\qquad c_a(G_t)=0\quad(t\in I).
                                                               \tag{C.76.1}
\]
The first prescription includes an interval without finite
endpoints. Let $Z\subset\overline D$ consist of the original
boundary circle, the actual lines over $A$, and all retained
strips. This is a closed subset, preserved by the loop and
by every actual interior operator wherever that operator
is defined. The restricted loop has a based contraction
through homeomorphisms of $Z$, with jointly continuous
inverses, fixing the original circle and commuting with
every $q\in Q$ on $Z\cap O$.

On every line over $a\in A$, the contraction is exactly
\[
                  (s,a)\longmapsto(s+r u_\tau(a),a),
                           \qquad 0\leq r\leq1.                \tag{C.76.2}
\]
On a strip retained by (C.76.1), choose C.75's actual
homeomorphism $b:I\to\mathbb R$ for one such endpoint.
The base contraction is
\[
 B_{r,\tau}(t)=b^{-1}\bigl(b(t)+r c_a(f_\tau)\bigr).
                                                               \tag{C.76.3}
\]
Both finite endpoint values, when there are two, are
prescribed simultaneously by (C.76.2). All original
component operators, including negative fiber multipliers,
are retained. One compact family of original closed-disk
homeomorphisms controls the forward and inverse evaluations
over all retained strips and all parameters. Hence the
construction includes infinitely many strips accumulating
at an interior line or at the original circle.

More precisely, let $\mathscr R_Z$ be the subgroup of the
original relative $Q$-centralizer preserving $Z$, and let
\[
 \begin{split}
 \mathscr A_Z={}&\{k\in\operatorname{Homeo}(Z):
          k|_{\partial\overline D}=\id,
          \ kq=qk\text{ on }Z\cap O\},\\
 \rho_Z:\mathscr R_Z&\longrightarrow\mathscr A_Z,
                         \qquad \rho_Z(h)=h|_Z.
 \end{split}                                                   \tag{C.76.4}
\]
These are the actual compact-open and inverse topologies,
and $\rho_Z$ is a continuous homomorphism. The contraction
just asserted fills the actual loop $\rho_Z(f_\tau)$ in
$\mathscr A_Z$. When every complementary interval is
retained, $Z=\overline D$ and this is a contraction in the
original relative centralizer itself. No lifting of the
constructed filling through $\rho_Z$ is asserted for a
proper $Z$.
\end{claim}

\begin{proof}
\emph{The endpoint image and its actual short increments.}
For any one strip with $k\in\{1,2\}$ finite endpoints,
write $c=(c_{a_1},\ldots,c_{a_k})$ and $V=c(G)\subset\mathbb R^k$.
This is the image of the actual group $G$, before taking
any closure. It is a real vector subspace. Indeed, every
$g\in G$ has the finite-word path (C.75.7) from the identity.
Its character image is a path from $0$ to $c(g)$ in $V$.
The planar three-word statement (C.37.3), using a zero
second coordinate for $k=1$, places every real fraction
$r c(g)$, $0\leq r\leq1$, in $V$. Addition and negation
then give all real multiples. This proves the assertion
and, in particular, that $V$ is closed in $\mathbb R^k$.
It is the real span of $\{c(f_\tau):0\leq\tau\leq1\}$:
one inclusion follows from generation, and the other from
the vector-space property just proved.

For $\delta>0$, retain the actual compact family
\[
 \begin{split}
 J_\delta&=\{f_s f_t^{-1}:s,t\in[0,1],\ |s-t|\leq\delta\},\\
 E_\delta&=J_\delta^6,
 \qquad D_\delta=c(J_\delta).
 \end{split}                                                   \tag{C.76.5}
\]
Both parameter families are compact and path connected,
contain the identity, and include all smaller choices of
$\delta$. Their global maps and inverses tend uniformly
to the identity as $\delta\downarrow0$, in the actual
closed-disk homeomorphism topology. The set $D_\delta$
is compact, symmetric and path connected. It spans $V$:
divide $[0,\tau]$ into intervals of length at most $\delta$
and telescope $c(f_\tau)$ into their character increments.

Choose a basis $v_1,\ldots,v_d$ of $V$ in $D_\delta$,
where $d\leq2$. For each $i$ a path in $D_\delta$ joins
$0$ to $v_i$. Again (C.37.3) gives
$[-v_i,v_i]\subset D_\delta+D_\delta+D_\delta$;
symmetry supplies the negative fractions. Consequently
\[
 \Bigl\{\sum_{i=1}^d a_i v_i:|a_i|\leq1\Bigr\}
       \subset c(E_\delta).
                                                               \tag{C.76.6}
\]
For $d<2$ use identity factors to reach six increments.
For $d=0$ the assertion is the identity image.
Thus $c(E_\delta)$ contains a neighborhood of $0$ in the
actual vector image $V$, for every $\delta>0$. This is a
statement about specified compact sets of actual source
words. It does not choose a continuous inverse to $c$.

Put $P=\{f_\tau\}\cup\{f_\tau^{-1}\}$. The same three-word
statement, now on the original prefix character path,
gives the uniform inclusion
\[
               r c(f_\tau)\in c(P^3)
                           \quad(0\leq r,\tau\leq1).
                                                               \tag{C.76.7}
\]
The single family $P^3$ works for every strip and both
endpoint characters, including varying ranks of $V$.

\emph{One compact family in the original relative group.}
Take the closure $L$ of $G$ in the group of actual
closed-disk homeomorphisms fixing the circle and commuting
with $Q$ on the interior. This is a Polish group: the
ambient compact-open homeomorphism topology controls maps
and inverses, and the stated pointwise equations are closed.
Every member of $L$ still has the principal form
\[
                  k(s,t)=(s+u_k(t),\beta_k(t)),
                                                               \tag{C.76.8}
\]
with $\beta_k$ increasing, fixing $A$ pointwise and
preserving every complementary interval. These are closed
conditions for maps and inverses. Each finite-endpoint
evaluation $c_a(k)=u_k(a)$ extends continuously to $L$.
For any particular endpoint vector its image on $L$ is
still $V$: density gives inclusion in the closed space
$V$, and the original image already equals $V$.
The translation equation (C.75.4) also passes to $L$ on
each retained strip by pointwise continuity. No assertion
about a quotient topology on $c(G)$ is used.

Choose a complete compatible metric on $L$. Exactly as
in the compact-product argument of C.63, choose
$\delta_n>0$ so that left multiplication by any member of
$E_{\delta_n}$ changes every previously formed product
by less than $2^{-n}$. This is possible by compactness
of the previous product family and the uniform identity
limit in (C.76.5). The families
\[
       K_n=E_{\delta_n}\cdots E_{\delta_1}P^3,
                       \qquad K_0=P^3,
                                                               \tag{C.76.9}
\]
are nested, and their infinite products converge uniformly
over the entire product of their compact parameter spaces.
Completeness gives a compact limit family $K\subset L$
containing every $K_n$. It is path connected: parameterize
$P$ by concatenating the original loop and its inverse,
and use the convex time bands in (C.76.5). The product
of these parameter spaces is path connected, and its
uniform product map into $L$ is continuous.
This includes convergence of the actual inverse maps.

Every member of $K$ fixes the original circle and commutes
with every labelled $q\in Q$. It preserves all principal
strips, whether retained or not. On a given strip the
functions $u_k(t),\beta_k(t)$ in (C.76.8) are jointly
continuous in its compact label and its real coordinate.

\emph{Base prescriptions and nonempty fiber intervals.}
On a strip with all finite-endpoint characters zero, use
C.27's based median contraction $B_{r,\tau}$ of the
original base loop. Its value is between $t$ and
$\beta_\tau(t)$, and hence is attained by an original
prefix value. Define
\[
 A_n(r,\tau,t)=\operatorname{conv}
       \{u_k(t):k\in K_n,\ \beta_k(t)=B_{r,\tau}(t)\}
                                                               \tag{C.76.10}
\]
on this strip. These are nonempty compact intervals.

On a strip satisfying (C.76.1), choose one of its eligible
endpoints $a$ and its actual coordinate $b$. Use (C.76.3),
and include \emph{all} finite endpoints in the vector $c$.
Define instead
\[
 A_n(r,\tau,t)=\operatorname{conv}
       \{u_k(t):k\in K_n,\ c(k)=r c(f_\tau)\}.
                                                               \tag{C.76.11}
\]
Equation (C.76.7) proves nonemptiness. For every label
occurring in (C.76.11), the $a$ coordinate and the exact
translation relation give
\[
                       \beta_k(t)=B_{r,\tau}(t).
                                                               \tag{C.76.12}
\]
Thus (C.76.11) really is an interval of displacements
in one specified principal fiber. Every one of its
labels has the prescribed displacement at \emph{both}
finite endpoints. The base map (C.76.3) lies between
$t$ and $\beta_\tau(t)$ even when $b$ is decreasing.
It is an increasing homeomorphism of $I$ with jointly
continuous inverse and all four based parameter faces.

In both definitions use $K$ in place of $K_n$ to obtain
a nonempty compact interval $A_\infty$. The four desired
fiber prescriptions are
\[
 \begin{split}
 U(0,\tau,t)&=0,\qquad U(1,\tau,t)=u_\tau(t),\\
 U(r,0,t)&=U(r,1,t)=0.
 \end{split}                                                   \tag{C.76.13}
\]
They belong to every $A_n$ on the respective faces.

\emph{The exact lower-limit comparison and relative selection.}
For each of (C.76.10) and (C.76.11) one has
\[
 A_n(p)\subset\liminf_{p'\to p}A_{n+1}(p'),
                     \qquad p=(r,\tau,t).
                                                               \tag{C.76.14}
\]
For (C.76.10), the short-increment argument of C.63 applies
directly. The base group is transitive on $I$ by C.75.
At every base point, arbitrarily short original increments
and their inverses move that point to both sides. The
inverse-evaluation argument (C.63.9) proves this assertion
for the actual specified point. These increments occur
in $E_{\delta_{n+1}}$, with five identity factors added.
The intermediate value theorem gives a correction of the
nearby target level, as close to the identity as desired.

For (C.76.11), take a label $k\in K_n$ with
$c(k)=r c(f_\tau)$. At a nearby parameter the required
character correction is the actual vector
\[
        v(p')=r'c(f_{\tau'})-c(k)\in V,
                              \qquad v(p')\longrightarrow0.
                                                               \tag{C.76.15}
\]
Choose any arbitrarily small $\delta'\leq\delta_{n+1}$.
By (C.76.6), every sufficiently nearby $p'$ has an actual
$e_{p'}\in E_{\delta'}$ with $c(e_{p'})=v(p')$.
Then $e_{p'}k\in K_{n+1}$ has precisely the new required
character vector. Equation (C.76.12) supplies its base
level at $t'$ automatically. As $\delta'\downarrow0$,
its fiber displacement at $t'$ tends uniformly to that
of $k$ near the chosen source point. This proves the
lower-limit inclusion for every unconvexified value.
Applying it to two endpoints and retaining their convex
coefficient proves (C.76.14) for the whole interval.
No continuous choice of $e_{p'}$ was used.

Replace the allowed sets by the prescribed singletons
(C.76.13) on the closed parameter boundary. The
lower-limit property survives by the same two comparisons
as in C.63: off that boundary it is (C.76.14), and along
the boundary it is continuity of the prescribed values.
Here are the precise selection ingredients. Choose at
each parameter a value of the $n$th interval, extend it
locally as a constant section, and use (C.76.14) to make
it arbitrarily close to the next interval. A locally
finite partition of unity combines these approximations,
since every fiber set is convex. Starting with error
$2^{-1}$, inductively choose continuous approximations
$U_n$ with distance less than $2^{-n}$ from the $n$th
modified interval and with
$|U_{n+1}-U_n|<2^{1-n}$. Their locally uniform limit is
continuous, belongs to the closed compact interval
$A_\infty$ at every parameter, and equals each prescribed
singleton exactly. This is C.63's relative selection
construction with the lower-limit hypothesis now proved
for the simultaneous character constraints.

\emph{Every original component operator.}
If $Q$ is connected, all its operators are the original
fiber translations and the selected maps commute with
them. Otherwise retain C.33's actual component generator
\[
                 q(s,t)=(\lambda s+\kappa,t+d),
                         \quad d>0,\quad\lambda\ne0.
                                                               \tag{C.76.16}
\]
All members of $L$, and hence all labels in $K$, obey
\[
       u_k(t+d)=\lambda u_k(t),\qquad
       \beta_k(t+d)=\beta_k(t)+d.
                                                               \tag{C.76.17}
\]
Consequently $A+d=A$, and the retained intervals are
permuted with all their exact endpoint criteria: the
character at $a+d$ is $\lambda c_a$.
For a proper interval $I\ne\mathbb R$, no nonzero power
of the base translation preserves $I$. Choose the
construction on one representative of each interval
orbit and transport it by the actual powers of $q$.
For the endpoint coordinate the transported choice is
$b'(t+d)=\lambda b(t)$, with any additive constant
permitted; (C.76.3) then transports exactly. The fiber
displacement transports by $U'(r,\tau,t+d)=\lambda U(r,\tau,t)$.
An operator stabilizing that proper interval belongs
to $Q^0$ and is an actual fiber translation, so this
transport is independent of its representative.

The only interval which can stabilize under a nonzero
component power is $I=\mathbb R$, when $A$ is empty.
It is retained by the zero-endpoint rule. Perform its
relative selection in the line bundle
$(t,z)\sim(t+d,\lambda z)$ exactly as in (C.63.20).
Its continuous fiber norm is $|\lambda|^{-t/d}|z|$;
the preceding approximation bounds are taken in this
norm. The convex fiber sets and their lower-limit
property are well-defined there by (C.76.17), and the
median commutes with the actual base translation.
This covers the negative values of $\lambda$ as well.
The constant $\kappa$ cancels in the commutation equation
and has not been removed from the operator label.

\emph{The homeomorphisms and their compact evaluation bounds.}
On each retained strip put
\[
 \begin{split}
 K_{r,\tau}(s,t)&=(s+U(r,\tau,t),B_{r,\tau}(t)),\\
 K_{r,\tau}^{-1}(s,v)&=
 \bigl(s-U(r,\tau,B_{r,\tau}^{-1}(v)),B_{r,\tau}^{-1}(v)\bigr).
 \end{split}
                                                               \tag{C.76.18}
\]
These are exact inverse homeomorphisms, jointly continuous
on each strip. Use (C.76.2) on lines over $A$ and the
identity on the original circle. The formulas preserve
each piece of $Z$ and have the four required based faces.
Their operator commutation follows from (C.76.16)--(C.76.17)
and the transported construction.

At every source point of a retained strip the selected
value lies between two $K$-values in one real fiber,
by (C.76.10)--(C.76.12). Connect their labels by a path
in the compact parameter space for $K$. The C.62 word
triangle, in precisely the form (C.63.14), gives
\[
       K_{r,\tau}(x)\in K K^{-1}K(x),\qquad
       K_{r,\tau}^{-1}(x)\in K^{-1}K K^{-1}(x).
                                                               \tag{C.76.19}
\]
For the inverse inclusion, the two increasing fiber
maps bracket the inverse point by their own inverses
in their common input fiber, as in (C.63.26).
On a line over $a\in A$, the real path
$\sigma\mapsto u_\sigma(a)$ joins $0$ to $u_\tau(a)$
on the prefix $[0,\tau]$. It attains $r u_\tau(a)$.
Thus the forward value in (C.76.2) is an original
prefix value, and its inverse is an original inverse
prefix value. Since $P\subset K$, the two bounds
(C.76.19) hold there as well. The bounds persist under
all component transports because each actual member
of $K$ commutes with $q$ on $O$.

\emph{All finite endpoint limits, including varying strips.}
Suppose $t_n\to a\in A$, with $t_n$ in retained strips
$I_n$, and allow both homotopy parameters to vary.
Choose a finite endpoint $a_n$ of $I_n$ lying between
$t_n$ and $a$. Then $a_n\to a$. For a zero-character
strip every $k\in K$ satisfies $u_k(a_n)=0$.
For a character-prescribed strip every label used
in its defining interval satisfies
\[
                     u_k(a_n)=r_n u_{\tau_n}(a_n),
                                                               \tag{C.76.20}
\]
because the vector constraint includes both finite
endpoints. This same equality holds for the zero-character
strip, whose right-hand side is zero.
Compactness of $K$ in the actual group makes
$u_k(t_n)-u_k(a_n)\to0$ uniformly in $k$: the principal
coordinate and its inverse are continuous on a compact
neighborhood of the finite line point under consideration.
Convexification preserves this estimate. It follows that
\[
         U(r_n,\tau_n,t_n)-r_n u_{\tau_n}(a_n)
                              \longrightarrow0.
                                                               \tag{C.76.21}
\]
The original compact loop makes its latter endpoint
values jointly continuous. Both base prescriptions
are between $t$ and $\beta_\tau(t)$, so their base
values tend to $a$ uniformly in the parameters.
Equations (C.76.20)--(C.76.21) give exactly (C.76.2)
at the limiting line, also when $a$ is not itself an
endpoint of one fixed interval.

For inverses, a base map lying pointwise between the
identity and $\beta_\tau$ has inverse lying pointwise
between the identity and $\beta_\tau^{-1}$. This follows
by applying the two increasing maps to the inequalities
at its inverse argument. Hence an inverse base argument
approaching $a$ also tends to $a$ uniformly. Applying
(C.76.21) at that argument proves the inverse limit
with displacement $-r u_\tau(a)$. Along $A$ itself,
both limits follow directly from (C.76.2).

\emph{The original circle and the exact restriction map.}
The two word families in (C.76.19) are compact families
of the original closed-disk maps, and every member
fixes the original circle pointwise. For $x_n\to\zeta$
on that circle, their evaluation values at $x_n$ tend
uniformly to $\zeta$. Thus (C.76.19) proves joint
continuity of both the constructed maps and their
inverses at every original boundary point, with no
restriction on the escaping principal coordinates
or on the varying strip index. The circle is fixed
throughout. The set $Z$ is closed because its complement
in the disk is exactly the union of the excluded open
strips. The preceding bijections and their continuous
inverses are therefore actual homeomorphisms of this
specified compact space, and the compact-parameter
evaluation criterion gives continuity into its actual
homeomorphism group.

Restriction to a closed invariant subset is continuous
for maps and inverses, so (C.76.4) is the asserted
continuous group homomorphism. Its source contains the
original loop. The formulas construct the claimed
filling of that loop's restriction; no assertion that
every intermediate map extends through an excluded
strip is made. The excluded intervals have a nonzero
finite-endpoint character but none satisfying (C.76.1).
C.75 identifies their nonzero endpoint-stabilizer images
as the full real line. Choosing a continuous,
boundary-controlled lift there remains a separate step
in C.12. In particular no source-word surjectivity alone
has been substituted for that lift.
\end{proof}
